\vspace*{40mm}
\begin{center}
{\Large\bfseries ZUR\par\medskip HANNOVERSCHEN TRIANGULATION.}
\end{center}
\clearpage\null\clearpage
\begin{center}
[I.]\par\medskip
{\Large\bfseries BRIEFWECHSEL}\par\medskip
{\bfseries ZUR HANNOVERSCHEN TRIANGULATION.}\par\bigskip
[1.]\par\medskip
\textsc{Gauss} an \textsc{Schumacher}.\quad Göttingen, 5.~Julius 1816.
\end{center}

\par\noindent\hspace*{\parindent}Vor allen Dingen meinen herzlichen Glückwunsch zu der herrlichen grossen Unternehmung, welche Sie mir in Ihrem letzten Briefe ankündigen. Diese Gradmessung in den k.~dänischen Staaten wird uns, an sich schon, über die Gestalt der Erde schöne Aufschlüsse geben. Ich zweifle indessen gar nicht, dass es in Zukunft möglich zu machen sein wird, Ihre Messungen durch das Königreich Hannover südlich fortzusetzen. In diesem Augenblicke kann ich zwar einen solchen Wunsch in Hannover noch nicht in Anregung bringen, da erst die Astronomie selbst noch so grosser Unterstützung bedarf: allein ich bin überzeugt, dass demnächst unsere Regierung, die auch die Wissenschaften gern unterstützt, dem glorreichen Beispiele Ihres trefflichen Königs folgen werde. Wir würden dann schon einen respectabeln Meridianbogen von \(6\frac{1}{2}\) Grad haben, und leicht würden sich dann auch noch diese Operationen mit den bayerischen Dreiecken in Verbindung setzen lassen. Letztere sind gewiss mit grösster Sorgfalt gemessen, und es ist zu beklagen, dass sie der Publicität entzogen werden.

\par\noindent\hspace*{\parindent}Über die Art, die gemessenen Dreiecke im Calcül zu behandeln, habe ich mir eine eigene Methode entworfen, die aber für einen Brief viel zu weitläufig sein würde. In Zukunft, falls ich bis dahin, wo Sie Ihre Dreiecke gemessen haben, sie nicht schon öffentlich bekannt gemacht haben sollte, werde ich mit Ihnen darüber umständlich conferiren: ja ich erbiete mich, die Berechnung der Hauptdreiecke selbst auf mich zu nehmen.
\par\medskip\noindent IX.\hfill44
\clearpage
\noindent346\hfill{\small\textbf{BRIEFWECHSEL MIT SCHUMACHER.}}\par\medskip

\par\noindent\hspace*{\parindent}Bei dem zweiten Theile Ihrer Unternehmung, der Messung des Längengrades, habe ich nur einen kleinen Zweifel. Ich meinte nemlich, dass die Länder der dänischen Monarchie eher \textls[120]{flach} zu nennen sind, wenigstens keine hohe Berge haben. Ist diese Voraussetzung gegründet, und sind Sie dann dadurch genöthigt, zur Bestimmung des astronomischen Längenunterschiedes einen Zwischenpunkt oder gar \textls[120]{mehrere} zu nehmen, so wird jener Bestimmung, auch wenn sie noch so geschickte Gehülfen und Hülfsmittel haben, doch immer eine kleine Ungewissheit ankleben.

\par\noindent\hspace*{\parindent}Einen grossen Vortheil haben Sie in dem Umstande, dass Dänemark schon einmal trigonometrisch vermessen ist; ich meine natürlich nicht in den gemessenen Winkeln selbst, die weit davon entfernt sind, sich zu einer Gradmessung zu qualificiren, sondern weil jene Operationen Ihnen das Auswählen der Stationspunkte ungemein erleichtern werden. Dies Aufsuchen würde mir bei einer ähnlichen Arbeit gerade das Unangenehmste sein, weil dabei so viele Zeit umsonst verloren wird. Ich habe mir viele Mühe gegeben (in ähnlichen Rücksichten auf künftige Operationen), die von \textsc{Epailly} im Hannoverschen gemessenen Winkel zu erhalten, aber ohne Erfolg.

\par\noindent\hspace*{\parindent}.\ .\ .\ .\ .\ .\quad Mir war eine interessante Aufgabe eingefallen, nemlich:
\begin{quote}
\noindent»\textbf{allgemein} eine gegebene Fläche so auf einer andern (gegebenen) zu projiciren (abzubilden), dass das Bild dem Original in den kleinsten Theilen ähnlich werde.«
\end{quote}
\par\noindent\hspace*{\parindent}Ein specieller Fall ist, wenn die erste Fläche eine Kugel, die zweite eine Ebene ist. Hier sind die stereographische und die merkatorsche Projection particuläre Auflösungen. Man will aber die \textls[120]{allgemeine} Auflösung, worunter alle particulären begriffen sind, für jede Art von Flächen.

\par\noindent\hspace*{\parindent}Es soll hierüber in dem Journal philomathique bereits von \textsc{Monge} und \textsc{Poinsot} gearbeitet sein (wie \textsc{Burckhardt} an \textsc{Lindenau} geschrieben hat), allein da ich nicht genau weiss, \textls[120]{wo}, so habe ich noch nicht nachsuchen können, und weiss daher nicht, ob jener Herren Auflösungen ganz meiner Idee entsprechen und die Sache erschöpfen. Im entgegengesetzten Fall schiene mir dies einmal eine schickliche Preisfrage für eine Societät zu sein.\quad.\ .\ .\ .\ .\ .
